\documentclass{article}
\usepackage{amsmath}
\usepackage{amsthm}

\newtheorem{theorem}{Problem}

\begin{document}

\begin{theorem}
Compute: $0.\overline{7}-0.\overline{4}+0.\overline{2}$. Express your answer as a common fraction.
\end{theorem}

\begin{proof}
Step 0: Step 1. I need to compute \( 0.\overline{7} - 0.\overline{4} + 0.\overline{2} \) and express the answer as a common fraction.
Step 1: Step 2. First, I should convert each repeating decimal into a fraction.
Step 2: Step 3. Let's start with \( 0.\overline{7} \). Let \( x = 0.\overline{7} \). Then, \( 10x = 7.\overline{7} \). Subtracting the original \( x \) from this, we get \( 10x - x = 7.\overline{7} - 0.\overline{7} \), which simplifies to \( 9x = 7 \), so \( x = \frac{7}{9} \). Therefore, \( 0.\overline{7} = \frac{7}{9} \).
Step 3: Step 4. Next, \( 0.\overline{4} \). Let \( y = 0.\overline{4} \). Then, \( 10y = 4.\overline{4} \). Subtracting, \( 10y - y = 4.\overline{4} - 0.\overline{4} \), which gives \( 9y = 4 \), so \( y = \frac{4}{9} \). Therefore, \( 0.\overline{4} = \frac{4}{9} \).
Step 4: Step 5. Now, \( 0.\overline{2} \). Let \( z = 0.\overline{2} \). Then, \( 10z = 2.\overline{2} \). Subtracting, \( 10z - z = 2.\overline{2} - 0.\overline{2} \), which gives \( 9z = 2 \), so \( z = \frac{2}{9} \). Therefore, \( 0.\overline{2} = \frac{2}{9} \).
Step 5: Step 6. Now, substitute these fractions back into the original expression: \( \frac{7}{9} - \frac{4}{9} + \frac{2}{9} \).
Step 6: Step 7. Since all fractions have the same denominator, I can combine them directly: \( \frac{7 - 4 + 2}{9} = \frac{5}{9} \).
Step 7: Step 8. Therefore, the answer is \( \frac{5}{9} \).

# Final Answer

\[ \boxed{\dfrac{5}{9}} \]
\end{proof}

\end{document}
